\begin{table}[htbp]
    \centering
    \caption{Recursos tecnológicos empleados en el proyecto.}
    \label{tab:recursos_tecnologicos}
    \renewcommand{\arraystretch}{1.30}
    \begin{tabularx}{\textwidth}{L L L C}
        \toprule
        \textbf{Recurso tecnológico} & \textbf{Detalle / Versión} & \textbf{Función en el proyecto} & \textbf{Costo / Licencia} \\
        \midrule
        Servicio de internet & Banda ancha 2 mensualidades & Investigación, repositorios y despliegue de la plataforma & 140 Bs/mes \\
        Servidor web & Apache 2.4 & Aloja la aplicación y atiende las solicitudes de inscripción & Open Source (0 Bs) \\
        Lenguaje de programación & PHP 8.2 & Lógica de negocio, sesiones y validación de credenciales & Open Source (0 Bs) \\
        Base de datos & PostgreSQL 16.4 & Guarda los datos y cuida los cupos sin duplicar & Open Source (0 Bs) \\
        Frontend & HTML5, CSS3, JavaScript ES6 & Interfaz adaptable y validación instantánea en el navegador & Open Source (0 Bs) \\
        Editor de código & Visual Studio Code & Escritura y organización del código fuente & Gratuito (0 Bs) \\
        Control de versiones & Git & Respaldo y control de cambios del proyecto & Gratuito (0 Bs) \\
        Navegadores web & Chrome / Firefox & Pruebas de compatibilidad en PC y celular & Gratuito (0 Bs) \\
        \bottomrule
    \end{tabularx}
    \notatabla{Entorno de desarrollo con herramientas libres, sin pago de licencias. El hardware y la papelería se detallan como recursos materiales en la Tabla~\ref{tab:recursos_materiales_det}.}
\end{table}
